\documentclass[11pt]{article}
\usepackage[utf8]{inputenc}
\usepackage[margin=1in]{geometry}
\usepackage{graphicx}
\usepackage{booktabs}
\usepackage{tabularx}
\usepackage{array}
\usepackage{enumitem}
\usepackage{float}

\setlength{\parindent}{0pt}
\setlength{\parskip}{0.5em}

\newcommand{\sampleframe}[1]{%
\begin{minipage}[t]{0.235\textwidth}
    \centering
    \includegraphics[width=\linewidth]{../../outputs/results/overlays/#1}\\
    {\scriptsize \texttt{\detokenize{#1}}}
\end{minipage}%
}

\begin{document}

\begin{center}
    {\Large \textbf{PitchVision Intermediate Report}} \\
    {\large 50\% Checkpoint: Visual Qualitative Results} \\
    \vspace{0.15cm}
    \textbf{Project:} Broadcast to 2D Mini-Map Projection \\
    \textbf{Author:} Peter Yaacoub \\
    \textbf{Subset:} 50 static frames with YOLO detections and team color labels
\end{center}

\vspace{0.2cm}
\hrule

\section{Checkpoint Goal}
This report documents the midpoint milestone defined in the project specification: the pipeline must successfully process single static frames and produce qualitative outputs with bounding boxes and team color labels. The current implementation covers player detection, jersey-color clustering, and overlay generation on the 50-image subset.

\section{Subset Summary}
The following aggregate results were produced from the 50-frame subset.

\begin{tabularx}{\textwidth}{@{} >{\raggedright\arraybackslash}p{0.45\textwidth} X @{} }
\toprule
\textbf{Metric} & \textbf{Value} \\
\midrule
Images processed & 50 \\
Total person detections & 536 \\
Average person detections per frame & 10.72 \\
Total ball detections & 38 \\
Average field lines detected per frame & 2.40 \\
Team accuracy & Not available yet, because the subset is currently used for qualitative validation \\
\bottomrule
\end{tabularx}

\section{Qualitative Results}
The overlays below show representative frames from the 50-image subset. Each output includes the detected bounding boxes and the team label assigned from jersey-color clustering. The selection is spread across the subset so the report captures different camera compositions, player densities, and lighting conditions.

\begin{figure}[H]
    \centering
    \sampleframe{frame_0001.jpg}\hfill
    \sampleframe{frame_0005.jpg}\hfill
    \sampleframe{frame_0010.jpg}\hfill
    \sampleframe{frame_0015.jpg}

    \vspace{0.35cm}

    \sampleframe{frame_0020.jpg}\hfill
    \sampleframe{frame_0025.jpg}\hfill
    \sampleframe{frame_0030.jpg}\hfill
    \sampleframe{frame_0035.jpg}

    \vspace{0.35cm}

    \sampleframe{frame_0040.jpg}\hfill
    \sampleframe{frame_0045.jpg}\hfill
    \sampleframe{frame_0048.jpg}\hfill
    \sampleframe{frame_0050.jpg}
    \caption{Representative qualitative outputs from the 50-image subset.}
\end{figure}

\section{Observed Behavior}
The visual outputs are consistent with the checkpoint objective:

\begin{itemize}[leftmargin=*]
    \item The detector reliably finds the majority of visible players in the broadcast frames.
    \item Team clustering is more stable on the regenerated run because the appearance features focus more tightly on the upper torso.
    \item Referee labels now appear only on a few central detections, which is preferable to over-assigning referee status to regular players.
    \item Ball detections are still sparse, but the lower confidence threshold and larger inference size recover the ball more often than before.
    \item Line extraction now returns a small, stable set of pitch markings again after restoring the baseline Hough settings.
\end{itemize}

\section{Checkpoint Status and Next Steps}
The 50\% milestone is complete for the static-image stage. The next implementation steps are to add explicit evaluation against manual labels, improve robustness of the team-color assignment on difficult crops, and then extend the pipeline to frame-to-frame tracking and pitch projection.

\end{document}